\section{JMP: treatment targets, selected subgroups and competing events}
\label{sec:exec24-jmp-reuse}

A bank analyst considering a retention campaign may want an overall
treatment effect, a subgroup worth targeting, or a prediction of default
before another event such as prepayment. JMP's retained documentation
provides distinct relevant workflows. Comparing them requires preserving
the estimand and edition rather than treating a causal-treatment report,
an uplift tree and a survival display as interchangeable evidence. The
proposal retains JMP as a commercial reference with conditional analyst
and worker routes; no native capability or entitlement is established
\cite{exec13_ch_jmp_causal,exec13_ch_jmp_uplift,exec14_survival_jmp,
exec24_jmp_choice}.

\subsection{Declared editions and the inspected technical boundary}
The causal-treatment pages are retained from the JMP Pro 19.2 help family
and display 13 January 2026 publication dates. The uplift overview belongs
to 19.0 and is dated 15 July 2025; the more detailed retained Consumer
Research guide is version 16. The supplied-model instrumental-variable
discussion is from the September 2025 version-19 Multivariate Methods
guide, while the selected survival pages are 18.1 and dated 27 June 2024.
Later acquisition of those mutable pages is not proof of exact bytes at
the 7 October 2026 cutoff. Older-guide behavior is not automatically
verified for a newer installed product
\cite{exec13_ch_jmp_causal,exec13_ch_jmp_uplift,exec13_ivrd_jmp,
exec14_survival_jmp}.

The causal help contains equations rendered as images. The retained
inspection uses the surrounding substantive prose, not an invented
transcription of those equations. The finite equations and examples below
are independently authored and identified as such. No additional image
review is claimed. Exact installed edition, supported automation profile,
commercial rights and prices remain separate questions; a help page is
neither a license nor evidence of an executed service
\cite{exec24_jmp_choice}.

\subsection{An average causal effect has a specified target population}
The documented Causal Treatment personality is available in JMP Pro and
supports binary or continuous treatment and outcome. For binary treatment,
the prose distinguishes average treatment effect from average treatment
effect among the treated. It documents IPWR, regression adjustment and
augmented inverse-probability methods using specified logistic propensity
and linear or logistic response models. For continuous treatment, the
reported targets are mean potential outcomes. These are useful concrete
capabilities to compare, but their interpretation depends on positivity,
consistency and conditional exchangeability
\cite{exec13_ch_jmp_causal}.

The assumption page explicitly distinguishes assumptions about the
data-generating process from empirical checks on a sample. A propensity
plot can reveal observed overlap problems; it cannot establish that every
unmeasured cause was included or that a campaign code corresponds to the
same intervention for every account. A business-to-causal assistant should
therefore record treatment versions, assignment information and alternative
causal explanations alongside any generated model specification. A fitted
response surface is not a substitute for that argument
\cite{exec13_ch_jmp_causal,exec24_jmp_choice}.

Targeting also requires preserving treatment-level coding. The version-16
uplift guide says that, with more than two treatment levels, the first
level is compared with a combination of all remaining levels. It also
documents value-order conventions and validation-column roles. These are
important historical contracts to test in the intended version, not
defaults that can be silently carried forward
\cite{exec13_ch_jmp_uplift}.

An authored randomized example has success probabilities $1/5$, $3/5$
and $9/10$ for a baseline and two distinct offers. If the pooled alternative
uses the latter offers with weights $(3/4,1/4)$, its contrast with baseline
is $19/40$. Reversing the mixture weights gives $5/8$. For 1,000
hypothetical accounts worth 100 USD per incremental success, the values
are 47,500 and 62,500 USD before offer costs. Both are valid mixture
contrasts under the stated experiment, but neither is automatically the
effect of a particular offer. The adapter must retain the first level,
pooled levels, their weights and the business action actually proposed
\cite{exec24_jmp_choice}.

\subsection{A selected uplift split needs its own evaluation law}
The retained uplift overview describes partition splits based on
treatment-by-split interactions, using linear models for continuous
responses and logistic models for categorical responses. It documents a
Monte Carlo adjustment intended to address predictor-selection bias from
different numbers of levels or candidate models. The older guide also
describes validation-based split selection. These features should be
acknowledged as documented controls, without inferring a general
post-selection interval or dependent-data guarantee
\cite{exec13_ch_jmp_uplift}.

A simple authored selection example explains the unresolved statistical
question. Suppose four independent candidate scores each equal minus one
or plus one with equal probability and have true mean zero. The selected
maximum is minus one only when all four scores are negative, which has
probability $1/16$. Its expected value is therefore
\[
 E[\max_j S_j]=\frac{15}{16}-\frac1{16}=\frac78.
\]
Selecting the largest score creates optimism even though no candidate
has a positive mean. This calculation does not test JMP's adjustment or
claim that it produces this bias. It specifies why the proposed comparison
needs a held-out or otherwise justified evaluation law for the complete
selection process, including predictor search, tree depth and policy choice
\cite{exec24_jmp_choice}.

In a bank, splits may be selected repeatedly across branches, customer
histories or calendar periods. Public fixtures should separate training,
selection and evaluation units according to the declared dependence regime.
They should retain all attempted splits and policies, treatment allocation,
costs and failure counts. A language-model agent may explain candidate
segments, but should not convert a selected in-sample difference into a
qualified intervention recommendation without the corresponding evidence
\cite{exec24_jmp_choice}.

\subsection{A continuous density model does not create empirical support}
The continuous-treatment prose specifies a normal conditional treatment
density with parameters estimated through linear regression. It describes
a stabilized-weight marginal structural model and a generalized propensity
score response model. These are particular modeling choices, not evidence
of identification for every dose or treatment distribution. A density
model assigning a positive number outside the observed range can conceal
an unsupported policy extrapolation if the wrapper records only that number
\cite{exec13_ch_jmp_causal,exec24_jmp_choice}.

For an authored bounded outcome, suppose treatments are observed only on
$[0,1]$. Two response laws agree there at $\mu(a)=a/2$. Beyond one, let
the first continue to $\mu_1(2)=1$, while the second declines according to
$\mu_2(a)=1-a/2$ and reaches $\mu_2(2)=0$. Both laws fit the same
observed support but imply opposite responses at the proposed dose two.
A parametric continuation chooses a restriction; it does not reveal the
missing response from the observed sample. Public tests can check that a
worker reports the support boundary and the chosen extrapolation model.
Actual off-support bank response requires new information or a justified
restriction \cite{exec24_jmp_choice}.

\subsection{Competing-event displays are not automatically removal interventions}
The inspected competing-causes page describes estimating each cause's
survival distribution by treating other events as censored, under a stated
independent-causes premise. It offers an Omit Causes display that
recalculates survival estimates to illustrate a situation in which selected
causes are no longer hazardous. Those documented operations do not establish
what would happen if a bank changed a prepayment policy or prevented a
particular competing event
\cite{exec14_survival_jmp}.

An authored finite example clarifies the distinction. Half of accounts
have a competing event at time one and latent default at time two; the
other half default at time two before their competing event at time three.
The observed fraction defaulting before a competing event is one half.
Under a specifically defined intervention that removes the competitor
while leaving each latent default time unchanged, all accounts default
by time two. Other intervention mechanisms could also change default
times. The observed competing-event record alone does not select the
unchanged-latent-time intervention. This example is not a native JMP
calculation and does not assume that its display uses this discrete model
\cite{exec24_jmp_choice}.

The proposed product therefore needs distinct labels for observed
cumulative incidence, cause-specific or net survival, and an explicitly
defined counterfactual removal target. Event coding, censoring, risk sets,
time units and dependence assumptions must be preserved. A plotting or
simulation option is a useful analytical interface; it does not supply the
causal intervention mechanism by itself
\cite{exec14_survival_jmp,exec24_jmp_choice}.

\subsection{The supplied-model IV route and remaining scope}
The retained version-19 multivariate guide documents model-implied
instrumental-variable estimation with equation-level diagnostics and
instrument lists. Its instruments follow a supplied structural model;
those declarations do not make invalid exclusion assumptions true or
establish unrestricted graph discovery. This is another conditional
comparison route, separate from causal-treatment regression and uplift
partitioning \cite{exec13_ivrd_jmp}.

The present evidence does not establish complete panel, sequential-treatment,
latent-discovery or agent-service contracts for every JMP product and
extension. That is a coverage limit, not a claim that the vendor lacks all
such features. The next survey should identify the intended edition,
official API or automation interface, supported artifact and entitlement,
then acquire the concrete missing technical pages. Image-only equations
need a fresh permitted visual inspection before claims relying on them
are admitted \cite{exec24_jmp_choice}.

Reference use is retained for the documented analytical workflows.
Conditional integration may use an entitled analyst workflow or isolated
worker after public target, coding, support, selection and event-time
fixtures pass. Direct adoption as the full causal product is deferred.
Porting a proprietary implementation is not an authorized open-source
reuse route; independent implementation of public methods is a separate
proposal. Unsupported causal, off-support or post-selection guarantees
are rejected without rejecting every bounded use of the software
\cite{exec24_jmp_choice}.

Scientific ownership covers intervention, target, assumptions and loss;
statistical ownership covers selection and dependence; adapter ownership
covers treatment order, mixture weights, events and typed refusals; platform
ownership covers supported automation, rights and complete costs. A
hypothetical 21,000 USD over 700 submitted, 525 completed and 350
qualified requests gives 30, 40 and 60 USD respectively. Actual prices,
support and bank value remain unmeasured. Native qualification, independent
review, visual and final acceptance remain open
\cite{exec24_jmp_choice}.
